\chapter{基于射线追踪的城市信道建模与深度射频指纹定位}

\section{引言}

% TODO: 本章研究动机
% - 城市环境中Ku/Ka波段信号的复杂传播特性（多径、NLOS、阴影衰落）
% - 传统统计信道模型无法刻画特定位置的信号空间结构
% - 需要高保真确定性信道建模 + 深度特征提取
% - 本章内容概述

\section{系统模型与问题描述}

\subsection{无人机辅助定位系统场景}

% TODO: 描述无人机搭载通信设备对地面非法终端定位的场景

\subsection{Ku/Ka波段信号传播模型}

% TODO: 卫星-地面链路的信号传播模型，包含直射、反射、绕射路径

\subsection{问题定义}

% TODO: 数学化定义定位问题——如何从多径畸变的接收信号中提取可靠的位置特征

\section{基于射线追踪的高保真三维城市信道建模}

\subsection{数字孪生仿真环境构建}

% TODO:
% - MATLAB Antenna Toolbox + RF Propagation Toolbox
% - OpenStreetMap导出北京市典型街区真实地理信息（建筑物高度、材质、几何轮廓）
% - 厘米级精度的三维电磁仿真环境

\subsection{Starlink基带信号生成与信道注入}

% TODO:
% - MATLAB 5G Toolbox + Phased Array System Toolbox
% - Starlink Ku波段帧结构信号生成（PSS/SSS/CSS）
% - 发射端模拟卫星多普勒频移，接收端模拟地面终端天线方向图
% - 将信号注入三维城市信道模型，捕捉时变CIR

\subsection{大规模电磁指纹数据集生成}

% TODO:
% - 覆盖目标区域27,546个网格点
% - 数据集内容：接收功率、到达角、时延扩展、I/Q双路复数信号
% - 数据集划分：训练集/验证集/测试集

\section{基于ResNet-50的深度射频指纹特征提取}

\subsection{I/Q双路信号预处理}

% TODO: 原始I/Q采样序列的预处理方法

\subsection{网络结构设计}

% TODO:
% - ResNet-50层级设计：{64, 64×3, 128×4, 256×6, 512×3}
% - 残差块结构与跳跃连接
% - 输出层设计（分类/回归）
% - 网络结构图

\subsection{训练策略与超参数配置}

% TODO: 学习率、批大小、优化器、损失函数、早停策略等

\section{MIMO空间分集增强}

\subsection{天线阵列配置方案}

% TODO: SISO → 2×2 → 4×4 → 8×8 MIMO的递进配置

\subsection{空间分集增益分析}

% TODO: MIMO如何利用空间分集增益增强信号特征区分度的理论分析

\section{实验结果与分析}

\subsection{实验环境与参数设置}

% TODO: 仿真环境参数、训练平台配置（PyTorch框架）

\subsection{定位精度对比实验}

% TODO:
% - ResNet-50 vs 传统几何方法：8.3m vs 45.2m
% - 精度提升一个数量级
% - 对比图表

\subsection{分类性能分析}

% TODO:
% - ResNet-50 F1-Score: 0.85 vs 4层浅层CNN: 0.73
% - 混淆矩阵分析

\subsection{MIMO规模对定位精度的影响}

% TODO:
% - 不同SNR条件下（-15dB、-10dB、-5dB、0dB），天线规模对精度的影响曲线
% - 8×8 MIMO在-15dB下仍控制在20m以内
% - 图：误差随天线规模变化趋势

\subsection{计算效率分析}

% TODO:
% - ResNet-50单次推理耗时约85ms
% - vs ResNet-101: 142ms
% - 精度与实时性的最佳平衡

\section{本章小结}

% TODO: 总结本章工作：
% 1. 构建了基于射线追踪的高保真城市信道模型
% 2. 提出了基于ResNet-50的深度RF指纹定位方法
% 3. 引入MIMO技术增强低SNR下的定位性能
% 4. 为下一章的多模态融合提供了高质量的单模态定位输入
